\section{Discussion}\label{sec:discussion}

The two empirical sections of this paper, the seven-model demonstration of \Cref{sec:demonstration} and the 30-paper re-analysis of \Cref{sec:phase2}, point in the same direction. Single-number ranking tables in published volatility-forecasting comparisons are systematically over-interpreted: when the differences are subjected to formal significance testing, the modal outcome is that the headline ``win'' is not statistically distinguishable from the runner-up.

\subsection{What the two empirical sections agree on}

The self-contained demonstration shows that within a single, well-documented seven-model panel, only 6 of 21 pairwise DM comparisons are significant at 5\%, the 90\% MCS keeps six of seven models, and SPA cannot reject any of the top six as the best forecaster. The 30-paper re-analysis shows that, across a larger sample, 86\% (12 of 14) of paper-claimed wins do not survive significance testing in our reproductions.

The agreement matters because the two pieces of evidence have orthogonal weaknesses. The seven-model demonstration is on one panel; its specific p-values are panel-specific. The Phase 2 reproductions span multiple asset classes, multiple model classes, and multiple authors but are necessarily approximations of the originals (RDS-1 papers ship no code; we re-implement from prose). Either piece of evidence taken alone would invite ``but is your specific test contaminated?'' as a response. Taken together, they describe a pattern that is robust to either critique.

\subsection{For practitioners deploying volatility forecasts}

The practical takeaway is that the choice between top-tier models in a typical volatility-forecasting comparison is often statistically indistinguishable. A risk manager comparing a GJR-GARCH against a LightGBM against an equal-weight ensemble on S\&P 500 data is not, in expectation, choosing between forecasters of different accuracy; she is choosing between forecasters of statistically equal accuracy with different operational properties (parameter count, retraining cost, interpretability, regulator-defensibility under SR 11-7 / OCC 2011-12 / EU AI Act). The honest framing of the choice is therefore not ``which is most accurate?'' but ``which is most defensible?'' \citep{FedSR117, OCC2011, EUAIAct}.

A second practical takeaway, less obvious from the original paper but visible in our \Cref{tab:demo_dm} and reinforced by the Phase 2 reproductions: the differences that \emph{are} statistically significant tend to be those that involve a clearly inferior model (HAR-RV losing to anything with a leverage-effect feature; GARCH losing to GJR-GARCH because the asymmetry term matters in this test period; in the Phase 2 sample, RNN clearly beating GARCH on Bitcoin in \citet{Shen2021} and ANN clearly beating GJR-GARCH on Bitcoin in \citet{MostafaSahaIslamNguyen2021}). The differences \emph{between} top-tier models, where the published literature most often draws its fine-grained ``X beats Y'' conclusions, are not significant. The implication is that an ensemble has a robustness argument even when it does not have an accuracy argument: it averages the regime-dependent failures of its components.

\subsection{For the methodological literature}

The single strongest empirical finding of this paper is from \Cref{sec:phase2_sigtest}: of the 30 published comparison papers in our sample, 27 (90\%) declare a winner without applying any formal predictive-accuracy test. This is the gap that \texttt{vol-eval} exists to close. The barrier to reporting a Diebold-Mariano p-value alongside any pairwise claim was always low; with a packaged tool, it is now zero. The remaining barrier is editorial: reviewers and editors need to expect, and where appropriate require, that the prose claim ``model X beats model Y on QLIKE'' come with the corresponding p-value.

Three editorial recommendations follow.

\textit{Default to reporting MCS survivors instead of single rankings.} A 90\% MCS that keeps six of seven models is informative content. A ranking table with bold values implying a clear winner where no clear winner exists is misinformation, even when each individual cell is correct. The MCS framing makes the underlying ambiguity visible.

\textit{Report Diebold-Mariano significance for any pairwise comparison the prose makes.} If the prose says ``model X beats model Y on QLIKE,'' the prose should also say at what p-value. If the answer is ``not significantly,'' the prose should say so plainly. \texttt{vol-eval}'s \texttt{dm\_test} returns the p-value as a one-line addition to any existing comparison.

\textit{Be especially careful with multi-model comparisons.} The DM test is for two models; using it as a battery of pairwise comparisons across many models without multiplicity adjustment will inflate the family-wise error rate and produce illusory ``wins.'' The MCS handles this correctly. The SPA / Reality Check handle the related ``one benchmark vs.\ many competitors'' question.

A fourth recommendation, which the Phase 2 disclosure findings (\Cref{sec:phase2_rds}) make especially salient: include a labelled ``Data and Code Availability'' block in the manuscript itself. Of the 30 papers in our sample, only the MDPI cohort consistently includes such a block (it is required by MDPI editorial policy as of 2020). The \emph{Journal of Financial Econometrics} cohort, despite being the most methodologically careful subset of our sample, embeds data descriptions in the prose of the data section without a separable disclosure block. This makes mechanical RDS scoring harder and reproduction-attempt feasibility opaque to the reader. A short labelled block costs the author nothing and saves the reader an hour of search.

\subsection{Limitations}\label{sec:limitations}

Five limitations of the present paper deserve careful noting.

First, the demonstration of \Cref{sec:demonstration} is on a single seven-model panel. The findings about MCS survival, SPA non-rejection, and the failure of the ensemble versus GJR-GARCH ``win'' to be significant are about that specific panel and that specific test period. The same demonstration on a different asset class, a different forecast horizon, or a different test window would produce different specific numbers. The methodological conclusion (that significance testing matters) generalizes; the specific p-values do not.

Second, the Phase 2 reproductions are necessarily approximations of the original papers. Because RDS-1 papers ship no code, we re-implement the headline architectures from each paper's prose description with standard hyperparameter defaults (1-layer LSTM with 50 hidden units, 20-day lookback, 30 epochs, RBF kernel for SVR with library defaults, GARCH(1,1) refit every 20 days with manual recursion). Where a paper discloses a different specification we follow the paper, but in many cases the disclosure is incomplete. This is itself a finding about the state of disclosure in the subfield, but it limits the strength of any individual ``does not survive'' verdict.

Third, three of the Phase 2 reproductions required substituting a library or feature we could not reproduce: \citet{Kim2021} needed PyMC for full Bayesian SV via MCMC (we approximated with AR(1) on log-squared returns), \citet{ErsinBildirici2023} needed the GARCH-MIDAS framework with macroeconomic indicators (we omitted both), and \citet{Kumar2024} needed PyTorch Geometric for the Temporal Graph Attention Network (we approximated with a multi-channel LSTM). For these three papers, the ``fails'' verdict in the headline tabulation is heavily caveated; we are not testing the paper's exact contribution, only a simpler counterpart. We retain these in the headline 86\% number for reporting honesty but flag them in \Cref{sec:phase2_caveats} and the per-paper \texttt{reproduction\_log.json} files.

Fourth, the 14-paper reproducible subset is not a random sample of vol-forecasting papers. It is the subset for which (a) the data is publicly accessible and (b) the headline architecture is reconstructable from the paper's prose. Papers requiring proprietary high-frequency tick data (the entire \emph{Journal of Financial Econometrics} cohort in our sample, six papers) are absent from the reproducibility analysis, and the absence is correlated with venue and methodology. The 86\% failure rate is the rate among the openly-reproducible subset, not the rate one would expect across the full 30-paper population.

Fifth, \texttt{vol-eval}'s implementation choices for the bootstrap and for the SPA centering follow \citet{Hansen2005} and the \texttt{arch} package's conventions. Alternative bootstrap designs (block bootstrap with fixed block length, parametric bootstrap, GARCH-type bootstrap) would produce slightly different p-values. We default to the stationary bootstrap because it is the standard in the published literature, but users who want a different design can pass it as an argument. The package's API does not foreclose alternatives; the current release simply ships one default.
